\subsection{Система метрик оцінювання якості перекладу}
\label{subsec:metrics}

Якість перекладу оцінено автоматичними метриками та ручною експертною
оцінкою. Автоматичні метрики вимірюють формальну відповідність референсному
тексту, тоді як експертна оцінка враховує семантичну адекватність та
природність мовлення.

\subsubsection{Автоматичні метрики}

\paragraph{BLEU (Bilingual Evaluation Understudy).}
Метрика BLEU~\cite{papineni_bleu} є однією з найбільш поширених метрик для
оцінки машинного перекладу. Вона базується на порівнянні $n$-грам
(послідовностей із $n$ слів) між машинним перекладом та одним або кількома
референсними перекладами.

Формула обчислення BLEU:
\begin{equation}
\text{BLEU} = BP \cdot \exp\left(\sum_{n=1}^{N} w_n \log p_n\right)
\end{equation}
де $p_n$ — модифікована точність (precision) для $n$-грам, $w_n$ — вагові
коефіцієнти (зазвичай $w_n = 1/N$), $N$ — максимальний порядок $n$-грам
(стандартно $N = 4$), а $BP$ — штраф за стислість (brevity penalty):
\begin{equation}
BP = \begin{cases}
1 & \text{якщо } c > r \\
e^{(1 - r/c)} & \text{якщо } c \leq r
\end{cases}
\end{equation}
де $c$ — довжина машинного перекладу, $r$ — ефективна референсна довжина.

Модифікована точність для $n$-грам обчислюється як:
\begin{equation}
p_n = \frac{\sum_{\text{C} \in \{\text{Candidates}\}} \sum_{\text{n-gram} \in C} Count_{clip}(\text{n-gram})}{\sum_{\text{C} \in \{\text{Candidates}\}} \sum_{\text{n-gram} \in C} Count(\text{n-gram})}
\end{equation}

BLEU обрано як стандарт де-факто для порівняння систем машинного
перекладу~\cite{post_bleu}: метрика дає швидке та відтворюване оцінювання
і корелює з людською оцінкою на великих корпусах.

\paragraph{BERTScore.}
На відміну від BLEU, метрика BERTScore~\cite{bert_score} оцінює семантичну
подібність на рівні контекстуалізованих векторних представлень слів,
отриманих за допомогою попередньо навченої мовної моделі BERT.

Для кожного токена $x_i$ в машинному перекладі та токена $\hat{x}_j$ в
референсі обчислюється косинусна подібність їхніх контекстуалізованих
ембедингів:
\begin{equation}
\text{sim}(x_i, \hat{x}_j) = \frac{\mathbf{x}_i^\top \hat{\mathbf{x}}_j}{\|\mathbf{x}_i\| \|\hat{\mathbf{x}}_j\|}
\end{equation}

На основі цих подібностей обчислюються три компоненти:
\begin{align}
P_{BERT} &= \frac{1}{|x|} \sum_{x_i \in x} \max_{\hat{x}_j \in \hat{x}} \text{sim}(x_i, \hat{x}_j) \\
R_{BERT} &= \frac{1}{|\hat{x}|} \sum_{\hat{x}_j \in \hat{x}} \max_{x_i \in x} \text{sim}(x_i, \hat{x}_j) \\
F_{BERT} &= 2 \cdot \frac{P_{BERT} \cdot R_{BERT}}{P_{BERT} + R_{BERT}}
\end{align}
де $P_{BERT}$ — точність, $R_{BERT}$ — повнота, $F_{BERT}$ — F1-міра.

BERTScore обрано, оскільки метрика краще враховує семантичну еквівалентність
перефразувань та синонімів і має вищу кореляцію з людською оцінкою
порівняно з BLEU для оцінки окремих речень, що важливо для
сегментованого мовлення.

\paragraph{METEOR (Metric for Evaluation of Translation with Explicit ORdering).}
Метрика METEOR~\cite{banerjee_meteor} розроблена для усунення деяких недоліків
BLEU. Вона враховує не лише точний збіг слів, але й стемінг (зведення до
основи слова), синоніми (на основі WordNet) та парафрази.

METEOR обчислює гармонійне середнє між точністю та повнотою уніграм з
урахуванням штрафу за фрагментацію:
\begin{equation}
\text{METEOR} = F_{mean} \cdot (1 - Penalty)
\end{equation}
де
\begin{equation}
F_{mean} = \frac{10 \cdot P \cdot R}{R + 9P}
\end{equation}
\begin{equation}
Penalty = 0.5 \cdot \left(\frac{\#chunks}{\#unigrams\_matched}\right)^3
\end{equation}

Тут $P$ — точність, $R$ — повнота, $\#chunks$ — кількість неперервних
фрагментів зіставлених уніграм.

METEOR має кращу кореляцію з людською оцінкою на рівні речень порівняно
з BLEU, особливо для мов із багатою морфологією. Врахування синонімів та
стемінгу є важливим для оцінки перекладу публіцистичних текстів.

\subsubsection{Людська експертна оцінка}

Автоматичні метрики, попри свою ефективність, не завжди адекватно відображають
якість перекладу з точки зору людини-читача. Тому в методології передбачено
можливість експертної оцінки за двома критеріями. На практиці основну увагу
зосереджено на анотації помилок та оцінці когерентності (розділ~\ref{sec:results}), оскільки
ці методи надають детальнішу діагностичну інформацію, ніж загальні шкали
адекватності та плавності.

\paragraph{Адекватність та плавність.}
Для оцінки якості перекладу можуть застосовуватися шкали
адекватності (ступінь збереження змісту оригіналу) та плавності
(природність мовлення цільовою мовою) за 5-бальною шкалою.
У роботі основну увагу зосереджено на більш діагностичних
методах --- анотації помилок за типами та оцінці міжфразової
когерентності, що надають детальнішу інформацію для аналізу каскадного
ефекту.

\subsubsection{Класифікація помилок перекладу}

Для детального аналізу помилок застосовано базову класифікацію, описану
в підрозділі~\ref{subsec:error_class} (пропуски, спотворення, додавання), розширену додатковими
категоріями на основі фреймворку MQM~\cite{mqm} (Multidimensional Quality
Metrics):
\begin{itemize}[nosep]
    \item \textit{Помилки в іменах власних} --- неправильна транслітерація
    імен, географічних назв, абревіатур;
    \item \textit{Граматичні помилки} --- порушення граматичних норм
    цільової мови;
    \item \textit{Стилістичні помилки} --- неприродні конструкції,
    калькування з мови оригіналу.
\end{itemize}

Кожна помилка додатково класифікується за ступенем серйозності:
\begin{itemize}[nosep]
    \item \textbf{Critical} (критична) --- помилка, що спотворює основний зміст
    повідомлення або містить фактичну неточність;
    \item \textbf{Major} (суттєва) --- помилка, що ускладнює розуміння фрагменту;
    \item \textbf{Minor} (незначна) --- помилка, що не впливає на загальне
    розуміння тексту.
\end{itemize}
